% =============================================================================
%  Spurious Precision: Why Shift-Share Designs Overstate Monetary Policy
%                      Effects on Local Employment
%  Conclusion — draft of August 11, 2026
%  Author: Abdallah Dalis
%
%  Step 7 of 00_Resources/research-paper-writing-process.md.
%  NO NEW RESULTS. Every number below appears earlier in the paper.
% =============================================================================

\section{Conclusion}

A shift-share difference-in-differences design applied to county employment
returns a precise, robust and wrong-signed estimate of how monetary policy
transmits to local labor markets. This paper asked what that precision was made
of, and answered that it was made of the business cycle.

Two diagnostics support the conclusion, and they fail in different ways.

% RESOLVED August 12, 2026: still two of four, but leads 2 and 3, not 2 and 4.
% This sentence does not name the leads, so it needed only the count. Section 5.5
% does name them and was changed with it.
The design fails its own pre-trends test. Adding leads of the interaction, two
of four reject zero under both clustered and randomization inference, and the
verdict holds across sample windows and permutation schemes.
Employment in exposed counties was already moving before policy changed, so the
estimate cannot be read causally.

Replacing the realized funds rate change with an identified high-frequency
surprise does not move the estimate to a different place. It widens the interval
% UPDATED August 12, 2026 for the rebuilt exposure: full window 5.30-9.97,
% excluded 5.80-11.65. The full window is the headline. Note the excluded range
% still rounds to "six to twelve" — do not promote it.
by a factor of five to ten in the full window, and six to twelve excluding
2009 to 2012, until no horizon rejects zero in either. The two estimators cannot
be told apart at any horizon. The collapse in precision, and not the change in
sign, is the finding.

Behind both is a measurement fact. Over this sample the realized rate change
correlates with contemporaneous national employment growth at $0.470$, while the
identified surprise correlates at $-0.119$. The endogenous measure carries the
cycle inside it, and a design built on that measure reports cyclical variation
as though it were policy variation.

\subsection*{What this implies for practice}

The recommendation is deliberately small, because a recommendation that requires
a new estimator will not be followed.

Report the design under both shock measures, and compare them on the standard
errors rather than the coefficients. A large gap is diagnostic. It says the
reported precision came from variation the identification strategy was meant to
remove. This costs one column and no new data.

The wider lesson is about how precision is read. A tight confidence interval is
usually taken as evidence that a design is working. In a shift-share setting it
can equally be evidence that the design is absorbing the business cycle, and the
two cases look identical in a published table. They are told apart by changing
the shock, not by inspecting the interval.

\subsection*{Limitations and next steps}

Three things this paper does not do.

It does not recover the true employment response. The identified estimates are
too imprecise to establish a sign, and claiming one would repeat the error the
paper documents.

It does not test exposure-robust inference. \citet{adao2019} show that
conventional standard errors can be badly understated when a small number of
common sectoral shocks drive exposure. Two permutation schemes were run and
neither addresses that concern: free permutation destroys the spatial
correlation of exposure, and within-state permutation narrows the null
mechanically. The appropriate instrument is their variance estimator, and
applying it is the first item of future work.

It does not describe the shape of the event-study path, in either direction.
The post-period coefficients move across defensible constructions of the lagged
variables, and the pre-period leads, while stable, are not monotone. The paper
reports coefficients and rejects zero. It does not narrate a trajectory that the
estimates would not support.

Beyond those, the natural extension is to ask whether the result generalizes.
The diagnostic proposed here is not specific to monetary policy or to county
employment. Any shift-share design whose shock is a realized policy variable
rather than an identified one is exposed to the same problem, and the same
one-column check would detect it.
